% Generated draft: requires mathematical and source review.
\documentclass[11pt]{article}
\usepackage[margin=25mm]{geometry}
\usepackage{fontspec}
\setmainfont{Latin Modern Roman}
\usepackage{amsmath,amssymb,mathtools,unicode-math}
\setmathfont{Latin Modern Math}
\usepackage{enumitem,needspace,xurl,hyperref}
\setlength{\parindent}{0pt}
\setlength{\parskip}{5pt}
\setlength{\emergencystretch}{3em}
\newcommand{\N}{\mathbb{N}}
\newcommand{\Z}{\mathbb{Z}}
\newcommand{\Q}{\mathbb{Q}}
\newcommand{\R}{\mathbb{R}}
\newcommand{\C}{\mathbb{C}}
\newcommand{\F}{\mathbb{F}}
\newcommand{\A}{\mathbb{A}}
\newcommand{\PP}{\mathbb{P}}
\newcommand{\E}{\mathbb{E}}
\newcommand{\eps}{\varepsilon}
\newcommand{\dd}{\,\mathrm{d}}
\begin{document}
\section*{UnsolvedMath ID: 10400120}
\textbf{Source problem code:} catalogue entry
\subsection{Definitions and mathematical statement}
Let \(\tau_r^G(M)\in\C\) be the source's quantum 3-manifold invariant for fixed admissible group and level data \((G,r)\), with its specified normalization and complex root conventions. Conjecture 7.5 asserts that on the nonzero locus its complex absolute value depends only on the abstract fundamental group.

Precisely, the natural pairwise formulation is: for manifolds \(M,N\) in the intended domain,
\[
 \pi_1(M)\cong\pi_1(N),\qquad
 \tau_r^G(M)\ne0,\quad\tau_r^G(N)\ne0
 \quad\Longrightarrow\quad
 |\tau_r^G(M)|=|\tau_r^G(N)|.
\]
The same \(G,r\), normalization and embedding into \(\C\) are used on both sides. Equivalently, the nonzero moduli factor through the group isomorphism class. The assertion does not claim that the complex phase is determined, that all levels or all groups give the same value, or that vanishing itself is a fundamental-group invariant.

The excerpt does not define \(\tau_r^G\), the allowed \(G,r\), or the manifold and auxiliary-structure category. The conventional closed-oriented 3-manifold quantum-invariant setting is suggested by the source context but is not fully stated here. These definitions and normalizations must be recovered before a concrete comparison can be made.
\subsection{Short English statement}
For fixed quantum-invariant parameters, do two manifolds with isomorphic fundamental groups have equal absolute invariant values whenever both values are nonzero? The nonvanishing restriction does not itself compare zero with nonzero values.
\subsection{Sources}
\begin{itemize}
\item Local UnsolvedMath registry, permanent ID 10400120, source code catalogue entry; no proposer or proposal year supplied; updated 2026-08-06.
\item Ohtsuki, \emph{Problems on invariants of knots and 3-manifolds} (2002), Conjecture 7.5, PDF page 102: \url{https://msp.org/gtm/2002/04/gtm-2002-04-024s.pdf}. Registry access date 2026-07-29; PDF font-structure extraction. Provenance distinguishes 231 numbered bold headings. Present primary-literature status and redistribution terms require review.
\item The supplied summary repeats the fundamental-group dependence conjecture without an exact additional reference or verified resolution.
\end{itemize}
\end{document}
